% Direct-authority diplomatic French for printed p.131.
\SourcePage{131}

\LayerTarget{tohoku:def:additive-diagram-relations:p2}
\noindent $R_{ij}$, on dira que $D$ est $R$-commutatif si tous les $D(L)$
$(L\in R)$ sont nuls. Appelons \textit{schéma de diagramme avec relations de
commutation} un couple $(S,R)=\Sigma$ formé d'un schéma de diagramme et d'un
ensemble $R$ comme ci-dessus. Pour toute catégorie additive $C$, on peut alors
considérer la sous-catégorie $C^\Sigma$ de $C^S$ formée des schémas
$R$-commutatifs.

\LayerTarget{tohoku:prop:1.6.1}
\noindent\textsc{Proposition 1.6.1.} --- \textit{Soient $\Sigma$ un schéma de
diagramme avec relations de commutation, $C$ une catégorie additive. Alors la
catégorie $C^\Sigma$ est une catégorie additive, et si $C$ est une catégorie
avec produits directs infinis (resp. sommes directes infinies) il en est de
même de $C^\Sigma$. De plus, si $C$ satisfait à un des axiomes AB 1) à AB 6)
ou des axiomes duals AB 3*) à AB 6*), il en est de même de $C^\Sigma$.}

\LayerTarget{tohoku:prop:1.6.1:proof}
D'ailleurs, si $D,D'\in C^\Sigma$, et si $u$ est un morphisme de $D$ dans
$D'$, alors son noyau (resp. conoyau, resp. image, resp. coimage) est le
diagramme formé par les noyaux (resp. \ldots) des composantes $u_i$, les
morphismes dans ce diagramme (correspondants aux flèches du schéma) se
déduisant de ceux de $D$ (resp. $D'\ldots$) de la façon évidente par
restriction (resp. passage au quotient). On interprète de façon analogue la
somme directe ou le produit direct d'une famille de diagrammes. Les sous-trucs
$D'$ d'un diagramme $D$ s'identifient aux familles $((D'(i))$ de sous-trucs des
$D(i)$ telles que pour toute flèche $\varphi$ d'origine $i$ et extrémité $j$,
on ait $D(\varphi)\,D'(i)\subset D'(j)$ ; alors $D'(\varphi)$ se définit comme
le morphisme $D'(i)\longrightarrow D'(j)$ défini par $D(\varphi)$. Les trucs
quotients de $D$ se déterminent de façon duale.

\LayerTarget{tohoku:def:dual-diagram-scheme}
Si $S$ est un schéma de diagrammes, on appelle \textit{schéma dual} et on
désigne par $S^\circ$ le schéma ayant mêmes sommets et même ensemble de flèches
que $S$, mais l'origine et l'extrémité des flèches de $S$ étant intervertis. Si
on donne de plus un ensemble $R$ de relations de commutation pour $S$, on
gardera le même ensemble pour $S^\circ$. Avec cette convention, pour une
catégorie additive $C$, la catégorie duale de $C^\Sigma$ s'identifie à
$(C^\circ)^{\Sigma^\circ}$.

\LayerTarget{tohoku:def:canonical-diagram-extension}
Soient $C,C'$ deux catégories additives, $\Sigma$ un schéma de diagramme avec
relations de commutation. Pour tout foncteur $F$ de $C$ dans $C'$, on définit
de façon évidente le foncteur $F^\Sigma$ de $C^\Sigma$ dans $C'^\Sigma$, dit
\textit{prolongement canonique de $F$ aux diagrammes}. $F^\Sigma$ se comporte
formellement comme un foncteur par rapport à l'argument $F$, en particulier un
momorphisme fonctoriel $F\longrightarrow F'$ définit un homomorphisme
fonctoriel $F^\Sigma\longrightarrow F'^\Sigma$. Notons enfin qu'on a pour un
foncteur composé :

\[
  (GF)^\Sigma=G^\Sigma F^\Sigma,
\]

\noindent et que les propriétés d'exactitude d'un foncteur se conservent par
prolongement à une classe de diagrammes.

\LayerTarget{tohoku:section:1.7}
\noindent\textbf{1.7. Exemples de catégories définies par des schémas de
diagrammes.}

\LayerTarget{tohoku:example:1.7:a}
\noindent a) Prenons $I$ réduit à un élément, et l'ensemble des flèches vide.
Alors les relations de commutation sont de la forme $n_i e=0$, donc peuvent se
réduire à une unique relation $ne=0$. Alors $C^\Sigma$ est la sous-catégorie de
$C$ formé des objets annulés par l'entier $n$. Si $n=0$, c'est $C$ lui-même.

\LayerTarget{tohoku:example:1.7:b:p1}
\noindent b) Prenons $I$ quelconque, pas de flèches, pas de relations de
commutation. Alors $C^\Sigma$ s'identifie à la catégorie produit $C^I$. Si on
suppose données des

\TohokuContinuation{132}
